\scanpage{78}
\mspage{69}
% archive sheet number: 78

$X\in\operatorname{supp}^{\circ}_M(A)$, on en conclut (puisque $\operatorname{supp}^{\circ}$ est fermé pour $\overset{\circ}{\ll}$) $\operatorname{omb}^{\circ}(X)\subset\operatorname{supp}^{\circ}_M(A)$, OK.

Ceci dit, les deux applications précédentes sont inverses l'une de l'autre, et établissent des iso. réciproques d'ens. ordonnés, munis de l'opération $\unsure{\complement|C}\ldots$

C'est formel en termes de
\diagram{\begin{tikzcd}[ampersand replacement=\&,row sep=2em,column sep=4em]
 L \arrow[d,hook] \arrow[dr,"x\mapsto\{x\}"] \arrow[dr,phantom,"\text{comm.}"{description,pos=.6,xshift=-15pt,yshift=-10pt}] \& \\
 M \arrow[r] \& \unsure{\mathcal{P}|P}(L)
\end{tikzcd}}
application (ici $X\mapsto\operatorname{omb}^{\circ}(X)$), $\mathrel{\|}_L$ relation sym. antiréfl. sur $L$ d'où relation $\mid\circ\mid_M$ sur $M$ par
\[
 X\mid\circ\mid Y\Longleftrightarrow o(X)\mathrel{\|}o(Y)
\]
i.e. $\forall x\in o(X),y\in o(Y),\ x\mathrel{\|}_L y$.
% margin note written slanted
\margin{On peut généraliser au cas d'une relation entre deux \uline{ens.} \uline{distincts} $M,M'$ versus $L,L'$, avec deux diagrammes au \deleted{\unsure{un}} lieu d'un, \unsure{fibré|figé}.}

% Several diagonal strokes cancel the following sentence and formula.
\deleted{À toute partie\inserted{$A$} \deleted{\unsure{de}} $\operatorname{supp}^{\circ}$ de $L$, on \unsure{associe} une partie\inserted{$\operatorname{supp}^{\circ}$} de $M$, par}
\[
\deleted{\unsure{\varphi|\phi}(A)=\{X\in M\mid o(X)\subset A\}}
\]
et on trouve \deleted{\unsure{ainsi}}\inserted{alors} deux iso \deleted{\unsure{tr}} réciproques.
\diagram{\begin{tikzcd}[ampersand replacement=\&]
\Sigma_L \arrow[r,shift left=.6ex] \& \Sigma_M \arrow[l,shift left=.6ex]
\end{tikzcd}}
\hfill $\ldots$

Il faudra bien que je me tape l'énoncé général, avec la vérification !

\uline{Cor} Si $\mathrel{\|}_L$ est la relation $x\ne y$ (cas \unsure{quasi-irréductible}) alors
\[
\Sigma_M\simeq\mathcal{P}(L)
\]
La notion de support\inserted{interne} est limpide : $\operatorname{supp}^{\circ}A$ s'identifie à $\operatorname{omb}^{\circ}(A)$.
